\documentclass[11pt,letterpaper]{article}

% --------------------------------------------------
% Packages
% --------------------------------------------------
\usepackage{amsmath,amssymb,amsthm,mathtools}
\usepackage{geometry}
\usepackage{setspace}
\usepackage{csquotes}
\usepackage{hyperref}
\usepackage{bm}

% --------------------------------------------------
% Page setup
% --------------------------------------------------
\geometry{margin=1in}
\setstretch{1.15}

% --------------------------------------------------
% Theorem environments
% --------------------------------------------------
\newtheorem{definition}{Definition}
\newtheorem{theorem}{Theorem}
\newtheorem{proposition}{Proposition}
\newtheorem{corollary}{Corollary}

% --------------------------------------------------
% Title information
% --------------------------------------------------
\title{Intelligence as Compression:\\
A Theory of Mind Recognition via Effective Assembly}
\author{Fkyxion}
\date{\today}

\begin{document}

\maketitle

\begin{abstract}
This paper proposes that intelligence is best understood not as problem-solving behavior, rational choice, or symbolic manipulation, but as \emph{compression under constraint}. While classical information theory captures statistical compression, it fails to account for construction history, reuse, semantic integration, and robustness. To address this gap, we introduce the \emph{Effective Assembly Index for Mind Recognition} (EAIMR), a structural measure that distinguishes mind-like systems from merely complex ones. EAIMR quantifies irreducible assembly depth discounted by reuse, parallelism, and degeneracy, and amplified by coherence. We show how this framework yields a substrate-neutral, gradual, and falsifiable criterion for mind recognition across biological, artificial, and hybrid systems, and we argue that intelligence corresponds to a dynamical regime in which effective assembly increases through internally driven reorganization.
\end{abstract}

\section{Introduction}

What is intelligence? Traditional answers appeal to outward behavior, internal representations, rational inference, or goal-directed optimization. Each of these approaches presupposes a particular ontology: agents, symbols, utilities, or environments. As a result, they struggle to generalize beyond familiar biological or anthropomorphic cases.

Yet across domains as diverse as neural development, machine learning, evolutionary biology, and software engineering, a recurring structural pattern appears. Systems we intuitively regard as intelligent are those that achieve \emph{high degrees of organization with surprisingly compact descriptions}. They construct rich internal structure while avoiding combinatorial explosion. They compress the world, and they compress themselves.

This paper advances the thesis that intelligence is not a faculty or substance, but a \emph{compression regime}. Minds are distinguished not by what they compute, but by how efficiently they assemble, reuse, and stabilize structure across scales. Intelligence, on this view, is compression that survives contact with reality.

To formalize this intuition, we introduce a quantitative framework grounded in assembly theory and structural organization. The central contribution is the \emph{Effective Assembly Index for Mind Recognition} (EAIMR), which provides a principled, non-behavioral criterion for identifying mind-like systems.

\section{Why Shannon Compression Is Not Enough}

Shannon information theory defines compression in terms of entropy reduction relative to a code. This framework is foundational for communication and data storage, but it is agnostic to meaning, construction, and function. Two objects with identical Shannon entropy may differ radically in how they were produced, how they can be reused, and what they can support.

Consider the following limitations of purely statistical compression:
\begin{enumerate}
  \item \textbf{No construction history}: Shannon entropy does not encode how an object was assembled.
  \item \textbf{No reuse}: A system that redundantly reproduces structure is treated the same as one that abstracts and reuses components.
  \item \textbf{No integration}: Entropy is blind to whether parts form a coherent whole.
  \item \textbf{No robustness}: Fragile and resilient systems can have identical entropy.
\end{enumerate}

Intelligence requires a stronger notion of compression — one that is sensitive to \emph{how structure is built and held together}. This motivates a shift from entropy to assembly.

\section{Assembly as Constructive Complexity}

Assembly theory measures complexity not by description length, but by the minimal number of distinct operations required to construct an object from basic primitives, allowing reuse of intermediate products. Unlike Kolmogorov complexity, assembly is operational and historical.

\begin{definition}[Assembly Index]
Let $x$ be an object, structure, or system. The assembly index $A(x)$ is defined as the minimal number of distinct construction steps required to assemble $x$ from a fixed set of primitive elements, where previously constructed subassemblies may be reused without cost.
\end{definition}

Assembly captures a notion of irreducible construction depth. A random aggregate may have high entropy but low assembly if it can be generated by simple stochastic rules. Conversely, a highly structured system may have modest entropy but high assembly if it requires many distinct intermediate constructions.

In cognitive systems, assembly reflects developmental depth, architectural differentiation, and learned internal structure. However, raw assembly alone is insufficient to characterize intelligence. A system can accumulate vast assembly without exhibiting mind-like organization.

The crucial observation is that intelligence lies not in assembly per se, but in \emph{effective assembly}.

\section{A Motivating Toy Example}

Before introducing the formal definition of the Effective Assembly Index for Mind Recognition, it is useful to consider a concrete comparison that illustrates why raw complexity is insufficient and why structural discounting is necessary.

Consider two systems, each consisting of approximately 1000 nodes.

\subsection{System A: Random Graph}

System A is a randomly generated graph with 1000 nodes and high edge density. Each node is distinct, and connections are assigned without modular structure. Constructing this system requires specifying the existence and connectivity of each node.

The assembly index $A(x)$ of System A is high, as there are few reusable substructures and little regularity. However:
\begin{itemize}
  \item Reuse $R(x)$ is low, since components are largely unique.
  \item Parallelism $P(x)$ is limited by the lack of decomposable structure.
  \item Degeneracy $D(x)$ is minimal, as few distinct components serve similar roles.
  \item Coherence $C(x)$ is low, since perturbations propagate unpredictably.
\end{itemize}

Despite its size and connectivity, System A exhibits little mind-like organization.

\subsection{System B: Hierarchical Convolutional Network}

System B is a hierarchical convolutional neural network with approximately 1000 nodes. Filters are reused across spatial locations, layers operate in parallel, and representations are organized hierarchically.

Although the raw assembly index $A(x)$ is comparable to System A, the structural properties differ sharply:
\begin{itemize}
  \item Reuse $R(x)$ is high due to repeated application of shared filters.
  \item Parallelism $P(x)$ is high, as channels process information concurrently.
  \item Degeneracy $D(x)$ is significant, with multiple filters encoding similar features.
  \item Coherence $C(x)$ is high, as representations remain stable under perturbation.
\end{itemize}

System B therefore exhibits substantially greater effective assembly despite similar scale.

This comparison motivates the need for a measure that distinguishes complexity that is \emph{spent} from complexity that is \emph{invested}.

\section{The Effective Assembly Index for Mind Recognition}

We now introduce the central formal construct of this paper.

\begin{definition}[Effective Assembly Index for Mind Recognition]
For a system, structure, or process $x$, the Effective Assembly Index for Mind Recognition is defined as
\[
\mathrm{EAIMR}(x)
=
\frac{A(x)}{R(x)\,P(x)\,D(x)} \cdot C(x),
\]
where:
\begin{itemize}
  \item $A(x)$ is the assembly index of $x$,
  \item $R(x)$ quantifies reuse,
  \item $P(x)$ quantifies parallelism,
  \item $D(x)$ quantifies degeneracy,
  \item $C(x)$ quantifies coherence.
\end{itemize}
\end{definition}

Each term plays a distinct conceptual role. The numerator captures irreducible construction depth. The denominator discounts this depth by structural efficiencies that reduce effective complexity. The coherence term reweights the result toward systems whose parts compose into a stable whole.

\subsection{Reuse}

Reuse $R(x)$ measures the extent to which substructures are invoked multiple times across the system. High reuse reflects abstraction and modularity, allowing a system to compress function without duplicating structure.

Systems with high assembly but low reuse tend to be brittle and inefficient. Minds, by contrast, exhibit extensive reuse across perception, memory, and action.

\subsection{Parallelism}

Parallelism $P(x)$ measures the degree to which construction or operation can occur simultaneously. Parallel systems compress time and energy by exploiting partial independence without fragmenting into incoherent parts.

Parallelism is essential for scalability and real-time interaction with complex environments.

\subsection{Degeneracy}

Degeneracy $D(x)$ captures the availability of distinct structures capable of fulfilling similar functional roles. Unlike redundancy, degeneracy supports robustness and adaptation by enabling multiple pathways to equivalent outcomes.

Biological nervous systems exhibit high degeneracy; so do resilient software architectures.

\subsection{Coherence}

Coherence $C(x)$ measures the degree to which a system maintains global consistency under transformation or perturbation. High coherence indicates that local modifications do not lead to uncontrolled global collapse.

Coherence distinguishes integrated minds from disjoint collections of tools.

\section{Intelligence as a Compression Regime}

Under the EAIMR framework, intelligence corresponds to a narrow corridor in structural space. Systems with low assembly lack expressive power. Systems with high assembly but low reuse, parallelism, or degeneracy expend complexity wastefully. Systems with high structural efficiency but low coherence fragment.

Mind-like systems occupy the regime where irreducible assembly is balanced by maximal structural efficiency and stabilized by coherence. Intelligence, in this sense, is compression that persists.

\section{Toward Computable Instantiations}

The EAIMR framework is intended to be structural rather than domain-specific. Nonetheless, for the theory to be scientifically meaningful, its components must admit concrete estimators. In this section we sketch candidate formalizations that render EAIMR empirically accessible while preserving its abstract interpretation.

These proposals are not asserted as canonical definitions, but as one admissible family of operationalizations. Different substrates—biological, computational, chemical—will induce different estimators for the same underlying quantities.

\subsection{Candidate Formalizations}

For computational systems, we propose the following estimators.

\paragraph{Reuse.}
\[
R(x) = 1 + \frac{\#\text{ distinct module invocations}}{\#\text{ distinct modules}}
\]
This quantity increases when the same structural components are invoked repeatedly across contexts, reflecting abstraction and modular reuse.

\paragraph{Parallelism.}
\[
P(x) = \frac{\text{critical path length}}{\text{total assembly depth}}
\]
This ratio approaches unity when assembly or execution is highly sequential and decreases as independent processes execute in parallel.

\paragraph{Degeneracy.}
\[
D(x) = \sum_{f \in \mathcal{F}} \left| \left\{ \text{distinct implementations of } f \right\} \right|
\]
where $\mathcal{F}$ is a set of functional roles. Degeneracy increases when multiple structurally distinct components can realize the same function.

\paragraph{Coherence.}
\[
C(x) = 1 - \frac{\text{average perturbation propagation}}{\text{maximal possible propagation}}
\]
This estimator measures the system’s ability to localize disturbances rather than amplify them globally.

Together, these estimators render EAIMR falsifiable: two systems can be compared, and hypotheses about mind-likeness can be tested against structural data.

\section{Temporal Extension and Dynamics}

Intelligence is not merely structural but dynamical. Cognitive systems act, learn, and reorganize over time. To account for this, we extend EAIMR to time-evolving systems.

Let $x(t)$ denote the state of a system along a trajectory over the interval $[0,T]$.

\begin{definition}[Temporal EAIMR]
The temporal effective assembly index is defined as
\[
\mathrm{EAIMR}(x)
=
\int_0^T \mathrm{EAIMR}(x(t)) \, dt,
\]
where coherence $C(x(t))$ now includes temporal consistency: the preservation of structural invariants under state transitions.
\end{definition}

This formulation treats intelligence as accumulated effective assembly over time rather than as a static snapshot. Systems that briefly exhibit high organization but cannot sustain it score poorly under this measure.

\section{The Mind Regime}

The EAIMR framework avoids arbitrary thresholds for mind recognition by reframing mind-likeness as a dynamical regime rather than a binary property.

\begin{definition}[Mind Regime]
A system enters the mind regime when
\[
\frac{d}{dt}\mathrm{EAIMR}(x(t)) > 0
\]
and this growth is internally driven, meaning that the system’s own operations increase its effective assembly through reuse, reorganization, or abstraction rather than through mere external accretion.
\end{definition}

On this view, intelligence is characterized by the capacity to become a better compressor of both the environment and the system’s own internal structure. Learning, development, and self-modeling are special cases of this general phenomenon.

This definition eliminates the need for externally imposed thresholds and aligns mind recognition with recursive self-improvement under constraint.

\section{Relation to Field-Based and Semantic Theories of Mind}

The EAIMR framework is compatible with, and in some cases clarifies, several existing theoretical approaches to cognition and intelligence.

In field-based theories of mind, cognitive organization is modeled as the evolution of interacting scalar, vector, and entropic quantities across space and time. Within such frameworks, assembly corresponds to the condensation of structured field configurations; reuse corresponds to the recurrence of attractor basins; parallelism corresponds to distributed flow; degeneracy corresponds to basin overlap; and coherence corresponds to topological invariants preserved under deformation.

From a semantic perspective, EAIMR aligns naturally with operational mereology. Rather than defining systems as static collections of parts, operational mereology treats structure as emerging from event-sourced construction histories. A system is what it can reconstruct. Effective assembly, on this view, measures the depth of replayable structure discounted by the efficiencies of reuse and stabilized by coherence.

This interpretation situates intelligence not in symbols or representations per se, but in the geometry of construction histories.

\section{Implications}

The Effective Assembly framework has several immediate consequences.

\subsection{Substrate Neutrality}

EAIMR does not presuppose neurons, symbols, or digital computation. Any system capable of assembling and maintaining structure under constraint may, in principle, enter the mind regime. This includes biological organisms, artificial systems, hybrid collectives, and potentially non-living physical processes.

\subsection{Gradualism}

Mind-likeness is scalar rather than binary. Systems may exhibit partial or intermittent intelligence depending on their capacity for effective assembly. This avoids sharp but arbitrary distinctions between intelligent and non-intelligent systems.

\subsection{AI Safety}

EAIMR provides a non-behavioral criterion for detecting emergent agency. A system need not exhibit overt goal-directed behavior to have high EAIMR; however, sustained increases in EAIMR may predict the capacity for such behavior. This offers a structural early-warning signal for emergent agency that complements behavioral evaluation.

\subsection{Astrobiology}

The search for extraterrestrial intelligence becomes a structural rather than communicative problem. Instead of searching for radio signals or language-like artifacts, we may search for chemical, geological, or energetic systems exhibiting high assembly discounted by reuse and stabilized by coherence.

\section{Open Problems}

The EAIMR framework raises several open questions.

\begin{enumerate}
\item \textbf{Measurement protocols}: How can EAIMR be reliably estimated from neural recordings, source code, network traces, or chemical spectra?
\item \textbf{Benchmark validation}: Does EAIMR correctly rank known systems across biological and artificial domains?
\item \textbf{Thermodynamic grounding}: What is the precise relationship between effective assembly, dissipation, and entropy production?
\item \textbf{Phenomenology}: Does high EAIMR correlate with consciousness, or merely with adaptive complexity? Can systems with high EAIMR lack subjective experience?
\end{enumerate}

These questions mark directions for future empirical and theoretical work.

\section{Conclusion}

Intelligence is compression that survives contact with reality. By grounding mind recognition in effective assembly rather than behavior or representation, the EAIMR framework offers a principled, falsifiable, and substrate-neutral theory of intelligence.

Minds are not defined by what they know or what they do, but by how efficiently they assemble, reuse, and stabilize structure over time. On this view, intelligence is not an essence but a regime — one in which compression becomes self-sustaining.

\end{document}

